\documentclass[12pt]{article}

% Page setup
\usepackage[a4paper, margin=1in]{geometry}

% Font
\usepackage{mathpazo}

% Formatting packages
\usepackage{setspace}
\usepackage{graphicx}
\usepackage{listings}
\usepackage{titlesec}
\usepackage{float} % For [H] placement

% Remove page numbers
\pagenumbering{gobble}

% Line spacing
\onehalfspacing

% Code style
\lstset{
    basicstyle=\ttfamily\small,
    frame=single,
    breaklines=true,
    tabsize=2
}

% Section formatting
\titleformat{\section}{\large\bfseries}{\thesection}{1em}{}

\begin{document}
\vspace{1cm}

\section*{Experiment Name}
Write a program to identify unique identifiers from a C program using a basic lexical analyzer.

\section*{Objective}
The objective of this experiment is to design and implement a basic lexical analyzer that scans a C program and detects valid identifiers. The program identifies unique identifiers while ignoring keywords, invalid tokens, and duplicates. This helps in understanding the working of lexical analysis in compiler design.

\section*{Algorithm}
\begin{enumerate}
    \item Start the program.
    \item Read the input C program from a file or user input.
    \item Define a list of C keywords (such as int, float, char, if, while, return, etc.).
    \item Initialize an empty set or list to store unique identifiers.
    \item Read the program line by line.
    \item For each line, break the line into tokens using delimiters such as spaces, commas, semicolons, parentheses, and operators.
    \item For each token:
    \begin{enumerate}
        \item Check if the token starts with a letter (a--z or A--Z) or an underscore (\_).
        \item Check if the remaining characters of the token contain only letters, digits, or underscores.
        \item If the token does not satisfy identifier rules, skip it.
        \item Compare the token with the list of C keywords.
        \item If the token matches a keyword, ignore it.
        \item Otherwise, check whether the identifier is already stored.
        \item If it is not present, add it to the list of unique identifiers.
    \end{enumerate}
    \item Continue scanning until the entire program is processed.
    \item Display all unique identifiers found in the C program.
    \item Stop the program.
\end{enumerate}

\newpage

\section*{Source Code}

\begin{figure}[H]
    \centering
    \includegraphics[width=0.7\textwidth]{source.png}
    \caption{Source code for lexical analyzer to detect unique identifiers}
\end{figure}
\begin{figure}[H]
    \centering
    \includegraphics[width=0.9\textwidth]{input.png}
    \caption{Input code for lexical analyzer to detect unique identifiers}
\end{figure}

\section*{Output Screenshot}

\begin{figure}[H]
    \centering
    \includegraphics[width=0.7\textwidth]{output.png}
    \caption{Output showing all unique identifiers extracted from the C program}
\end{figure}

\section*{Discussion}
This experiment demonstrates the design of a simple lexical analyzer to extract identifiers from a C program. The analyzer scans the input, applies identifier rules, removes keywords, and ensures uniqueness using a suitable data structure. This experiment provides a clear understanding of tokenization, identifier recognition, and the importance of lexical analysis in the first phase of a compiler.

\end{document}
